\section{Distinctive Predictions and Contact With Data (P26, P27, P35)}
\label{sec:predictions}

A model becomes science when its distinctive predictions meet data. Three of ours do.

\subsection{Swerves (P26)}

On a continuum geodesic a free particle keeps its velocity forever. On a causal set
there is no element exactly straight ahead, so the particle hops to the nearest future
element and its momentum gets a tiny random kick each step. \textbf{Where we landed:}
the rapidity diffuses, its variance growing linearly with proper time (a diffusion with
no continuum analogue, the Dowker-Henson-Sorkin swerve). The diffusion is present at
every sprinkling density. The honest caveats: the variance saturates at large proper
time (a finite-box effect) and the precise density scaling needs a covariant trajectory
rule. The swerve itself is robust.

\subsection{Dark energy, the headline match (P35)}

The everpresent $\Lambda$ predicts $\delta\Lambda \sim 1/\sqrt{V}$. At the observed
spacetime $4$-volume (the past horizon, $(ct/\ell_{\mathrm{P}})^4$):
\begin{itemize}
\item predicted $\delta\Lambda \sim 1.5\times10^{-122}$ (Planck units),
\item observed $\Lambda \sim 2.9\times10^{-122}$ (Planck units),
\item ratio $0.53$.
\end{itemize}
\textbf{Where we landed:} the prediction matches the measured dark energy to order of
magnitude, within a factor of two. The right value, about $10^{-122}$, comes out as one
over the square root of the spacetime volume, not as a tuned constant. This is the one
quantitative cosmological success of causal-set theory.

\subsection{Lorentz violation, a confirmed null (P27, P35)}

Gamma-ray-burst photon timing (GRB 090510, Fermi LAT) excludes linear Lorentz violation
below about $7.6$ times the Planck energy. The framework, using a random sprinkling that
is Lorentz-invariant in distribution, predicts \emph{no} linear Lorentz violation at
all. \textbf{Where we landed:} observation confirms the prediction, and it discriminates,
a lattice (or any regular) discreteness would predict a large effect that is already
excluded, while the random-sprinkling discreteness predicts the observed null. The
discreteness scale is bounded only from below.

\subsection{The swerve as a future target}

The swerve diffusion is Planck-suppressed, so the accumulated spread over the age of the
universe is small, below current bounds (from the narrowness of cosmic-ray and particle
momenta). It is a future target for ultra-high-energy cosmic rays rather than a present
conflict.

\textbf{Bottom line.} One prediction matches a measured number (dark energy, to order of
magnitude), one is confirmed by a null result that excludes the alternatives (no linear
Lorentz violation), and one is a clean future target (swerves). This is where the
framework makes contact with data, not just internal coherence.
